\chapter{Detecção e Monitoramento}
\label{cap:deteccao}

A capacidade de detectar atividades maliciosas em tempo real, sem degradar o
desempenho das aplicações monitoradas, é um dos principais diferenciais do eBPF
frente a soluções tradicionais de segurança \cite{cassagnes2020}.
Este capítulo apresenta técnicas concretas de detecção e monitoramento baseadas
em eBPF, organizadas pelos vetores de observação mais relevantes em ambientes de
produção.

% ----------------------------------------------------------------
\section{Detecção de execução de processos}
% ----------------------------------------------------------------

A execução de processos é um dos eventos mais informativos para segurança:
praticamente toda ação maliciosa — desde a exploração de uma vulnerabilidade até
a persistência no sistema — envolve em algum momento a criação ou substituição de
um processo \cite{gregg2019}.
O eBPF permite monitorar esse vetor com cobertura total e latência mínima,
independentemente da origem do processo.

O \textit{tracepoint} \texttt{sys\_enter\_execve} é acionado toda vez que um
processo invoca a chamada de sistema \texttt{execve(2)}, responsável por substituir
a imagem do processo atual por um novo executável \cite{linuxkernel2024}.
Um programa eBPF nesse ponto coleta, sem modificar o processo:

\begin{itemize}
  \item PID e PPID — identificadores do processo e de seu pai, permitindo
        reconstruir a árvore de execução;
  \item UID e GID — identidade real do solicitante, revelando tentativas de
        execução sob identidade indevida;
  \item caminho do executável e argumentos da linha de comando, via
        \texttt{bpf\_probe\_read\_user\_str};
  \item cgroup e namespace de PID, essenciais para identificar o contêiner de
        origem em ambientes Kubernetes \cite{cassagnes2020}.
\end{itemize}

O exemplo a seguir ilustra um script bpftrace que registra toda execução
originada de processos filhos de servidores web — padrão típico de exploração de
vulnerabilidades de execução remota de código:

\begin{lstlisting}[language=bash, caption={Detecção de shell a partir de processo de serviço web com bpftrace.}]
tracepoint:syscalls:sys_enter_execve
/comm == "nginx" || comm == "apache2"/
{
    printf("[ALERTA] shell filho de %s: %s (pid=%d uid=%d)\n",
           comm, str(args->filename), pid, uid);
}
\end{lstlisting}

Além do \texttt{execve}, o \textit{tracepoint} \texttt{sched\_process\_fork}
captura a criação de processos via \texttt{fork(2)} e \texttt{clone(2)},
permitindo rastrear a propagação de processos maliciosos que se replicam antes de
executar sua carga \cite{gregg2020}.
Combinados em um mapa BPF que mantém a árvore de parentesco em memória, esses
dois pontos de observação fornecem uma visão completa do ciclo de vida de processos
sem qualquer instrumentação das aplicações \cite{majors2022}.

% ----------------------------------------------------------------
\section{Monitoramento de arquivos, chamadas de sistema e conexões}
% ----------------------------------------------------------------

\textbf{Acesso a arquivos.}
Programas eBPF conectados ao \textit{tracepoint} \texttt{sys\_enter\_openat}
registram toda abertura de arquivo com contexto completo — processo, caminho,
\textit{flags} de acesso e resultado da operação \cite{linuxkernel2024}.
Caminhos de particular interesse para segurança incluem:

\begin{itemize}
  \item \texttt{/etc/shadow}, \texttt{/etc/passwd} e \texttt{/root/.ssh/} —
        arquivos de credenciais frequentemente alvo de exfiltração;
  \item \texttt{/proc/\$PID/mem} e \texttt{/proc/\$PID/maps} — interfaces de
        memória de processo usadas por técnicas de injeção \cite{manpages2024};
  \item \texttt{/usr/bin/}, \texttt{/lib/} e \texttt{/sbin/} abertos em modo de
        escrita — indicativo de \textit{rootkits} que modificam binários do sistema.
\end{itemize}

A ferramenta \texttt{opensnoop} do BCC \cite{bcc} implementa esse padrão e é
amplamente utilizada em investigações forenses para identificar quais arquivos
foram acessados durante um incidente.

\textbf{Chamadas de sistema.}
O conjunto completo de \textit{tracepoints} \texttt{sys\_enter\_*} oferece
visibilidade sobre qualquer chamada de sistema invocada por qualquer processo
\cite{linuxkernel2024}.
Para fins de detecção, chamadas críticas são monitoradas seletivamente para
reduzir o volume de eventos:

\begin{itemize}
  \item \texttt{ptrace} — interceptação de outro processo; abusada por
        \textit{malware} para injeção de código e por \textit{rootkits} para
        ocultar presença;
  \item \texttt{setuid} / \texttt{setgid} / \texttt{setcap} — alteração de
        identidade ou capacidades, indicativas de escalonamento de privilégios;
  \item \texttt{mmap} com \texttt{PROT\_EXEC} — mapeamento de memória executável,
        padrão de carregamento de shellcode \cite{gregg2019};
  \item \texttt{init\_module} / \texttt{finit\_module} — carregamento de módulos
        de kernel, vetor clássico de instalação de \textit{rootkits}.
\end{itemize}

\textbf{Conexões de rede.}
Combinando \textit{tracepoints} de \textit{socket} com \textit{hooks} de XDP e
TC, o eBPF oferece visibilidade em múltiplas camadas da pilha de rede
\cite{linuxkernel2024}.
O monitoramento de conexões de saída — via \texttt{sys\_enter\_connect} — permite
identificar processos que estabelecem comunicação com endereços não previstos,
padrão associado a \textit{beaconing} de \textit{Command and Control} (C2) e
exfiltração de dados \cite{cassagnes2020}.
A ferramenta \texttt{tcpconnect} do BCC \cite{bcc} registra cada nova conexão TCP
com PID, endereço de destino e porta, sem impacto perceptível sobre o throughput
da rede \cite{gregg2019}.

% ----------------------------------------------------------------
\section{Detecção de comportamento anômalo}
% ----------------------------------------------------------------

A detecção de anomalias vai além do monitoramento de eventos individuais: consiste
em identificar sequências ou padrões de comportamento que diferem do que é esperado
para um processo, serviço ou usuário \cite{majors2022}.
O eBPF é particularmente adequado para esse fim porque permite acumular estado em
mapas BPF ao longo do tempo e realizar correlações diretamente no kernel, antes
mesmo que os eventos sejam enviados ao espaço de usuário \cite{linuxkernel2024}.

\textbf{Desvio de perfil de chamadas de sistema.}
Cada classe de aplicação — servidor web, banco de dados, agente de monitoramento
— apresenta um perfil estável de chamadas de sistema em operação normal.
Um programa eBPF pode construir, durante um período de observação inicial, um
histograma de frequência de chamadas por processo e alertar quando a distribuição
observada diverge significativamente desse perfil de referência
\cite{gregg2019}.
Essa técnica é resistente a \textit{malware} que tenta mimetizar nomes de
processos legítimos, pois o comportamento em nível de chamadas de sistema é mais
difícil de falsificar.

\textbf{Comunicação lateral entre contêineres.}
Em ambientes Kubernetes, comunicação direta entre \textit{pods} que normalmente
não se comunicam é um indicador de movimentação lateral após comprometimento
\cite{cassagnes2020}.
Mapas BPF do tipo \texttt{BPF\_MAP\_TYPE\_HASH} indexados por par
(cgroup de origem, cgroup de destino) permitem detectar esses fluxos anômalos em
tempo real, sem depender de logs de \textit{firewall} com atraso de segundos
\cite{linuxkernel2024}.

\textbf{Escalada de privilégios em tempo de execução.}
A sequência \texttt{fork} → \texttt{execve} de um binário SUID → \texttt{setuid(0)}
é o padrão de exploração de vulnerabilidades de escalonamento de privilégios local
mais comum \cite{gregg2020}.
Um programa eBPF que rastreia transições de UID em mapas indexados por PID detecta
essa sequência em microssegundos, muito antes que ela se manifeste em logs de
auditoria \cite{gregg2019}.

\textbf{Injeção de código em processos em execução.}
O uso de \texttt{ptrace} ou de \texttt{process\_vm\_writev} para escrever em
memória de outro processo é uma técnica de injeção amplamente explorada
\cite{gregg2019}.
Correlacionando o PID do processo chamador com o PID alvo e verificando se a
relação pai-filho é legítima — usando a árvore de processos mantida em um mapa
BPF — é possível distinguir operações legítimas de depuradores de tentativas de
injeção maliciosa \cite{majors2022}.

% ----------------------------------------------------------------
\section{Auditoria de eventos}
% ----------------------------------------------------------------

A auditoria de segurança tem como objetivo registrar, de forma confiável e
verificável, eventos relevantes para conformidade e investigação forense
\cite{beyer2016}.
O subsistema \texttt{auditd} do Linux é a solução tradicional para esse fim, mas
apresenta limitações de desempenho em sistemas de alta carga e dificuldades de
extensão para novos tipos de eventos \cite{gregg2020}.
O eBPF surge como complemento ou substituto natural: coleta eventos com menor
sobrecarga, permite filtragem inteligente diretamente no kernel e é extensível sem
alterações no sistema operacional \cite{cassagnes2020}.

\textbf{Registro de eventos com contexto completo.}
Por meio do mapa \texttt{BPF\_MAP\_TYPE\_RINGBUF}, programas eBPF emitem registros
de auditoria estruturados para o espaço de usuário com campos como timestamp em
nanossegundos, PID, UID, cgroup, chamada de sistema, argumentos e código de retorno
\cite{linuxkernel2024}.
Essa riqueza de contexto supera a granularidade oferecida pelo \texttt{auditd}
para a maioria dos eventos \cite{gregg2019}.

\textbf{Filtragem inteligente no kernel.}
Diferentemente do \texttt{auditd}, que envia todos os eventos para o espaço de
usuário e filtra depois, programas eBPF aplicam regras de interesse diretamente
no kernel — registrando apenas eventos que satisfazem critérios de segurança
predefinidos \cite{linuxkernel2024}.
Isso reduz drasticamente o volume de dados transmitidos e processados, tornando
viável a auditoria contínua mesmo em sistemas com dezenas de milhares de chamadas
de sistema por segundo \cite{cassagnes2020}.

\textbf{Integridade do registro de auditoria.}
Uma preocupação legítima em auditoria forense é a possibilidade de um atacante com
privilégios suficientes suprimir registros de suas ações.
Programas eBPF em execução no kernel são mais difíceis de desabilitar
silenciosamente do que agentes em espaço de usuário: sua remoção requer privilégios
explícitos e gera eventos observáveis por outros programas eBPF — criando
uma cadeia de custódia parcialmente auto-referencial \cite{majors2022}.

\textbf{Integração com sistemas externos.}
Registros de auditoria coletados por programas eBPF podem ser encaminhados para
sistemas de gerenciamento de eventos e informações de segurança (SIEM) via
exportadores como o \texttt{ebpf\_exporter} \cite{ebpfexporter} ou via
\textit{pipelines} Prometheus/Grafana \cite{prometheusdocs,grafanadocs}.
Essa integração centraliza os dados de múltiplos nós em uma visão unificada,
essencial para detecção de ataques distribuídos que envolvem vários sistemas
simultaneamente \cite{beyer2016}.